\begin{table}[t]
\centering
\scriptsize
\caption{Primary sparse comparison on the 97 standard nonnegative instances. BW denotes backward weight; lower values are better. Best counts every method attaining the minimum observed BW on an instance; ties are credited to all tied methods. Relative excess is the mean percentage reduction achieved by IPSNS under the baseline-normalized convention \((X-\mathrm{IPSNS})/X\times 100\), where \(X\) is the comparator BW on completed instances.}
\label{tab:sparse-external-baselines}
\begin{tabular}{@{}lrrrrr@{}}
\toprule
Method & Complete & Mean BW & Mean RT (s) & Best & Relative excess \\
\midrule
IPSNS (ours) & 97/97 & 37,698 & 21.92 & 96 & 0.00\% \\
LR-TA (ours) & 97/97 & 38,327 & 0.08 & 80 & 0.71\% \\
WMSF seed & 97/97 & 40,005 & 1.31 & 61 & 2.06\% \\
DRMacIver/FAS & 93/97 & 53,173 & 4.00 & 56 & 21.61\% \\
igraph Eades & 97/97 & 95,920 & 0.01 & 40 & 30.54\% \\
Weighted Eades & 97/97 & 99,689 & 0.11 & 42 & 30.48\% \\
Borda net score & 97/97 & 512,277 & 0.00 & 27 & 55.55\% \\
Random multistart & 97/97 & 1,075,258 & 0.03 & 42 & 49.76\% \\
\bottomrule
\end{tabular}
\vspace{2pt}

{\footnotesize\raggedright DRMacIver/FAS is the main external comparator on this benchmark, but it completes 93 of 97 instances because two DAG cases return empty-tournament errors and two large sparse instances time out in the wrapper used here. Its mean relative excess of 21.61\% is a comparator-normalized reduction relative to IPSNS, not an IPSNS-normalized percent-worse score.}
\end{table}
